\lesson{1}{Mar 18 2022 Fri (23:11:52)}{Intercepts/Asymptotes/Domain/Range}{Unit 5}

\begin{definition}[Intercept]
  \label{def:intercept}

  When the line intersects/touches/goes through/etc the $x$ or $y$ axis. If the
  line touches the $x$-axis, then it's a $x$-intercept. Likewise, if the line
  touches the $y$-axis, then it's a $y$-intercept.
\end{definition}

\begin{definition}[Asymptotes]
  \label{def:asymptotes}

  A line that a curve approaches as it heads toward infinity.
\end{definition}

\begin{definition}[Domain]
  \label{def:domain}

  The set of all numbers that appear as $x$-coordinates of points on the graph.
\end{definition}

\begin{definition}[Range]
  \label{def:domain}

  The set of all numbers that appear as $y$-coordinates of points on the graph.
\end{definition}

\paragraph{Find Vertical/Horizontal asymptotes}
\label{par:find_vertical_horizontal_asymptotes}

\begin{marginfigure}
    \centering
    \incfig{vertical-and-horizontal-asymptotes}
    \caption{Vertical and Horizontal Asymptotes}
    \label{fig:vertical_and_horizontal_asymptotes}
\end{marginfigure}

\begin{itemize}
  \label{item:find_vertical_horizontal_asymptotes}

  \item Vertical Asymptote(s):

    To find the vertical asymptote, you try to find where the graph of the
    function are approaching the $y$-axis, but never touch. From this, we
    conclude that the vertical asymptote is $x = 1$.

  \item Horizontal Asymptote(s)

    To find the horizontal asymptote, you try to find where the graph of the
    function are approaching the $x$-axis, but never touch. From this, we
    conclude that the vertical asymptote is $y = 3$.
\end{itemize}

% paragraph find_vertical_horizontal_asymptotes (end)

\paragraph{Domain and Range}
\label{par:domain_and_range}

\begin{itemize}
  \label{item:domain_and_range}

  \item Domain:

    The domain is usually the set of all real numbers except the values of the
    vertical asymptotes. That's only if the graph doesn't have any holes. We
    found the vertical asymptote as $x = 1$. So, here's the domain:
    $(-\infty, 1) \cup (1, \infty)$ 

  \item Range:

    We project the graph onto the $y$-axis. The numbers that appear as
    $y$-values of points on the graph will make our range. Remember not to
    include the horizontal asymptote. So, here's our range:
    $(-\infty, 3) \cup (3, \infty)$
\end{itemize}

% paragraph domain_and_range (end)

\paragraph{Find all $x$-intercepts and $y$-intercepts}
\label{par:find_all_x_intercepts_and_y_intercepts}

\begin{itemize}
  \label{item:find_all_x_intercepts_and_y_intercepts}

  \item Find all $x$-intercepts:

    To find the $x$-intercepts, just look for when the graph passes through the
    $x$-axis. So, the $x$-intercept is $-2$.

  \item Find all $y$-intercepts:

    To find the $y$-intercepts, just look for when the graph passes through the
    $y$-axis. So, the $y$-intercept is $-6$.
\end{itemize}

% paragraph find_all_x_intercepts_and_y_intercepts (end)

\newpage
